% =====================================================================
% ejercicio03.tex
% Responsabilidad única: mostrar, con un ejemplo numérico y su
% generalización algebraica, por qué una variable de escala mayor
% domina la distancia euclidiana cuando no se estandariza — es la base
% conceptual que se retoma empíricamente en el Ejercicio 7.
% =====================================================================
\subsection{Ejercicio 3 --- Efecto de la escala sobre la distancia}

Se consideran los puntos $A=(20,1000)$, $B=(25,1200)$, $C=(40,1010)$.

\begin{enumerate}

\item \textbf{Distancias en escala original.}
\[
    d(A,B) = \sqrt{(20-25)^2+(1000-1200)^2} = \sqrt{25+40000} = \sqrt{40025} \approx 200.06
\]
\[
    d(A,C) = \sqrt{(20-40)^2+(1000-1010)^2} = \sqrt{400+100} = \sqrt{500} \approx 22.36
\]

\item \textbf{Vecino más cercano.}
Como $d(A,C)\approx 22.36$ es mucho menor que $d(A,B)\approx 200.06$,
el punto \textbf{C} es el vecino más cercano a $A$.

\item \textbf{Distancias tras dividir la segunda variable entre 100.}
Con $A'=(20,10)$, $B'=(25,12)$, $C'=(40,10.1)$:
\[
    d(A',B') = \sqrt{5^2+2^2} = \sqrt{29} \approx 5.385
\]
\[
    d(A',C') = \sqrt{20^2+0.1^2} = \sqrt{400.01} \approx 20.0002
\]

\item \textbf{¿Cambió el vecino más cercano?}
Sí, ahora $d(A',B')\approx 5.39 < d(A',C')\approx 20.00$ así que el
vecino más cercano pasó a ser \textbf{B}, el punto opuesto al de antes
de reescalar pero esto tiene una explicación directa: En la escala
original la segunda variable (diferencias en centenas) dominaba por
completo el cálculo, mientras que al dividirla entre 100 su diferencia
se vuelve comparable a la de la primera variable y esta última antes
irrelevante, ahora pasa a determinar la distancia.

\item \textbf{Por qué una escala mayor domina la distancia.}
En la expresión general
\[
    d(x,z) = \sqrt{\sum_{j=1}^{d} (x_j-z_j)^2},
\]
cada variable contribuye a través del cuadrado de su diferencia, de
modo que si la variable $j$ tiene un rango considerablemente mayor al
de las demás (como la segunda coordenada (ordenada) de este ejercicio en orden de miles y no decenas), $(x_j-z_j)^2$ será en términos absolutos,
mucho mayor que los demás términos, de forma que tanto la suma como,
en consecuencia, la distancia total terminan determinadas casi
exclusivamente por esa variable sin importar si conceptualmente es la
más relevante para el problema, este es precisamente el motivo por
el cual antes de aplicar un algoritmo basado en distancias como KNN,
se recomienda estandarizar o normalizar las variables.

\end{enumerate}

\begin{figure}[H]
\centering
\begin{tikzpicture}[every node/.style={font=\footnotesize}]

% ===== Panel 1: antes de estandarizar =====
\begin{scope}
\draw[gray!50] (0,0) -- (3.5,0);
\pgfmathsetmacro{\hA}{200.06*0.011}
\pgfmathsetmacro{\hC}{22.36*0.011}
\draw[fill=gray!45, draw=gray!70] (0.15,0) rectangle (0.85,\hA);
\draw[fill=classgreen, draw=black] (2.15,0) rectangle (2.85,\hC);
\node at (0.5,\hA+0.28) {$200.06$};
\node[classgreen, font=\scriptsize\bfseries] at (2.5,\hC+0.32) {$\star$ más cercano};
\node at (2.5,\hC+0.62) {$22.36$};
\node at (0.5,-0.35) {$d(A,B)$};
\node at (2.5,-0.35) {$d(A,C)$};
\node[align=center, font=\small\bfseries] at (1.5,-1.0) {Antes de estandarizar};
\end{scope}

% ===== Panel 2: despues de estandarizar =====
\begin{scope}[shift={(5.6,0)}]
\draw[gray!50] (0,0) -- (3.5,0);
\pgfmathsetmacro{\hBp}{5.385*0.11}
\pgfmathsetmacro{\hCp}{20.0002*0.11}
\draw[fill=classgreen, draw=black] (0.15,0) rectangle (0.85,\hBp);
\draw[fill=gray!45, draw=gray!70] (2.15,0) rectangle (2.85,\hCp);
\node[classgreen, font=\scriptsize\bfseries] at (0.5,\hBp+0.32) {$\star$ más cercano};
\node at (0.5,\hBp+0.62) {$5.385$};
\node at (2.5,\hCp+0.28) {$20.0002$};
\node at (0.5,-0.35) {$d(A',B')$};
\node at (2.5,-0.35) {$d(A',C')$};
\node[align=center, font=\small\bfseries] at (1.5,-1.0) {Después de estandarizar};
\end{scope}

\end{tikzpicture}
\caption{Cambio del vecino más cercano de $A$ tras dividir la segunda
variable entre 100 (cada panel usa su propia escala vertical, ya que
las magnitudes antes y después no son comparables): pasa de ser $C$ a
ser $B$.}
\end{figure}
